% !TEX root = ../main.tex
\documentclass[../main.tex]{subfiles}

\begin{document}

\chapter{Role Categories in Sustainability Transitions}
\labch{background}
\margintoc

\begin{chapteroverview}
Before we can map role constellations, we must understand how the field has conceptualized them. This chapter surveys role-based perspectives in sustainability transitions research, examines how major frameworks handle actor roles, and pays particular attention to intermediaries---the focal actors of this dissertation. We trace intermediary conceptualization from bilateral brokers to systemic actors, examine their functions and dynamics, and conclude by identifying six persistent problems that motivate our cartographic response. This directly addresses our first sub-research question on developing a conceptual framework for theory-guided yet empirically open role analysis.
\end{chapteroverview}


%===============================================================================
\section{From Actors to Roles to Constellations}[Actors to Constellations]
\labsec{background:intro}
%===============================================================================

Sustainability transitions are fundamentally \textbf{multi-actor processes}. No single organization---however powerful---can transform an energy system, reshape a food regime, or reconfigure urban mobility alone. Transitions emerge from the coordinated and contested actions of diverse actors: startups and incumbents, policymakers and activists, researchers and users, investors and intermediaries.

Early transition frameworks acknowledged this diversity but offered limited vocabulary for analyzing it \cite{markard2012sustainability}.\sidenote{\textbf{Theoretical gap}---Frameworks distinguish actor types but not how roles evolve and interrelate.} The Multi-Level Perspective distinguished niche innovators from regime incumbents. Strategic Niche Management identified network builders and expectation managers. Transition Management convened frontrunners in transition arenas. Yet these frameworks treated roles implicitly---as background assumptions rather than objects of analysis.

In recent years, scholars have argued that understanding \emph{who does what} in transitions---and how actors' roles evolve and combine---is crucial for explaining why some transitions accelerate while others stall \cite{wittmayer2017roles,fischer2016actors,avelino2017power}. This has generated growing interest in \textbf{role constellations} as analytic tools for unpacking social dynamics of systemic change.

\begin{roledef}[Role Constellation]
A \textbf{role constellation} is a configuration of actor roles and their interrelations within a transition process. It describes not merely which actors are present, but what functions they perform and how those functions combine to enable (or constrain) transition dynamics.
\end{roledef}

The formal introduction of role constellations to transition research is often traced to \textcite{wittmayer2017roles}, who imported sociological role theory to provide vocabulary for analyzing changing actor interactions over time. They argued that transition studies needed to move beyond identifying actor \emph{types} (who is present) to analyzing actor \emph{roles} (what functions actors perform and how these combine). This distinction matters: Multiple actors from different sectors might perform the \emph{same} role (several organizations acting as brokers). One actor can occupy \emph{multiple} roles simultaneously (an intermediary that convenes, funds, and advocates). And roles can \emph{shift} over time as transitions unfold (a niche protector becoming a regime challenger).

Crucially, roles are not fixed attributes but \textbf{emergent and malleable}. They can be created (establishing a new ``mobility broker'' function), altered (an incumbent shifting from resistor to adapter), or deliberately assigned (a governance intervention that empowers certain actors to perform coordination roles). This malleability makes role constellations not merely an analytical lens but a potential \textbf{intervention point} for transition governance.


%===============================================================================
\section{Roles Across Theoretical Frameworks}[Roles in Theory]
\labsec{background:frameworks}
%===============================================================================

How do the major transition frameworks conceptualize roles? Each offers a partial map with characteristic strengths and blind spots.\sidenote{\textbf{Different lenses}---Each framework illuminates some roles while obscuring others.} The \textbf{Multi-Level Perspective} distinguishes niche actors from regime incumbents \sidecite{geels2002technological}. \textbf{Transition Management} emphasizes frontrunners who experiment in protected arenas \sidecite{kemp2007tm}. \textbf{Strategic Niche Management} identifies cosmopolitan actors who aggregate lessons across local experiments \sidecite{smith2007snm}. \textbf{Innovation Systems} approaches specify seven system-building functions from knowledge development to legitimation \sidecite{bergek2008tis}. These frameworks overlap without explicit integration, each highlighting different roles within transition processes.


%===============================================================================
\section{Fischer \& Newig: A Systematic Review of Actor Typologies}[Fischer \& Newig]
\labsec{background:fischer-newig}
%===============================================================================

\textcite{fischernewig2016actors} provide the most comprehensive mapping\sidenote{\textbf{Fischer \& Newig (2016)}---Systematic review of 386 articles identifies four overlapping typologies.} of how sustainability transitions research categorizes actors. Analyzing 386 articles from 1995--2014, they addressed whether actors have been neglected in transitions literature in favor of abstract system concepts.

Their key finding: the field lacks a unified actor typology. Multiple overlapping classification schemes coexist:

\begin{itemize}[nosep]
\item \textbf{Systemic typology}: niche, regime, landscape actors---derived from MLP
\item \textbf{Institutional typology}: state, market, civil society---the governance triangle
\item \textbf{Governance typology}: local, regional, national, global scales
\item \textbf{Intermediary typology}: actors bridging other categories
\end{itemize}

These typologies coexist without integration. A startup can be simultaneously a niche actor (systemic), market actor (institutional), and local actor (governance). The literature conflates these without systematic treatment of how they relate.

\begin{insightbox}[The Overlap Problem]
A single organization might be categorized as a ``niche actor'' (systemic), ``civil society organization'' (institutional), ``local actor'' (governance), and ``intermediary'' (bridging function)---all simultaneously. Without explicit attention to how these framings intersect, research findings become difficult to cumulate across studies.
\end{insightbox}

A crucial finding across frameworks: actor roles are \textit{erratic}---they change over time, and actors can belong to different categories simultaneously. Most typologies assume static category membership, missing the temporal dynamics that characterize actual transitions. Fischer and Newig note that ``agency, strategic behaviour, and learning dynamics of actors are somewhat underdeveloped'' in the literature, despite their centrality to transition processes.

The review also reveals that intermediaries function as a cross-cutting category---they don't fit neatly into systemic, institutional, or governance typologies because they are defined by \textit{position} (in-between) rather than \textit{attributes}. This supports a relational framing where actor ``type'' matters less than relational configuration.


%===============================================================================
\section{Wittmayer et al.: Role Theory in Sustainability Transitions}
\labsec{background:wittmayer}
%===============================================================================

\textcite{wittmayer2017roles} made the seminal contribution\sidenote{\textbf{Wittmayer et al. (2017)}---Introduces sociological role theory to transition studies. This dissertation explicitly builds on and operationalizes their framework through multimodal cartography.} of explicitly importing sociological role theory into sustainability transitions research. Drawing on \textcite{biddle1986role}, they distinguish multiple conceptualizations of ``role'':

\textbf{Functional roles} emphasize positions within social systems and associated behavioral expectations. Organizations are expected to perform certain functions based on their structural position.

\textbf{Symbolic interactionist roles} emphasize meaning-making and identity construction. Roles are not just performed but negotiated through interaction.

\textbf{Organizational roles} focus on formal positions within hierarchies and the expectations attached to those positions.

\textbf{Cognitive roles} emphasize mental schemas and scripts that guide behavior in particular situations.

This multiplicity explains persistent definitional ambiguity: when researchers use ``role,'' they may invoke any of these meanings without explicit acknowledgment.

Beyond these definitional distinctions, Wittmayer et al. identify \textbf{three complementary perspectives} on how roles function in transitions:\sidenote{\textbf{Three perspectives}---Each perspective suggests different analytical strategies, which this dissertation operationalizes through different data modalities.}

\begin{enumerate}[nosep]
\item \textbf{Roles as recognizable activities}: What actors \emph{do}---observable behaviors and practices that can be identified and categorized
\item \textbf{Roles as resources}: Following \textcite{callero1994limits}, roles are ``cultural objects'' that actors strategically deploy to achieve goals
\item \textbf{Roles as boundary objects}: Following \textcite{simpson2008role}, roles serve as points of negotiation where different actors contest meaning and attribution
\end{enumerate}

These three perspectives map directly onto the multimodal cartography developed in this dissertation. \textbf{Websites} capture roles as recognizable activities---what organizations claim to do. \textbf{Social media} reveals roles as resources---how actors strategically deploy role claims in interaction. \textbf{News media and visual materials} expose roles as boundary objects---where external observers attribute, contest, and negotiate role meanings. This operationalization transforms Wittmayer et al.'s theoretical framework into systematic empirical method.

Crucially, Wittmayer et al. emphasize that roles are \textbf{dynamic and contested}.\sidenote{\textbf{Role dynamics}---Roles are actively constructed and negotiated, not merely occupied.} Actors actively construct roles through strategic action; they negotiate role boundaries with other actors; they resist role expectations that conflict with their interests or identities. This perspective transforms roles from static categories into \textbf{objects of struggle}.

The framework identifies three key role dynamics:

\begin{enumerate}[nosep]
\item \textbf{Role-taking}: adopting established roles with their associated expectations
\item \textbf{Role-making}: actively shaping and redefining role content
\item \textbf{Role creation}: establishing entirely new roles without precedent
\end{enumerate}

These dynamics are particularly important during transitions, when established role structures are disrupted and new roles emerge. The concept of \textbf{role constellation}---webs of roles that interact, interrelate, and co-evolve around specific issues---shifts analysis from individual roles to relational configurations. Each framing of a persistent problem carries implicit ideas about roles and responsibilities.

However, their operationalization remains qualitative and case-specific. How does one identify ``shared role understandings'' empirically? Their Rotterdam neighborhood case provides rich description but limited systematic method for tracking role constellation change at scale.


%===============================================================================
\section{De Haan \& Rotmans: Agency and Transformative Actors}[De Haan \& Rotmans]
\labsec{background:dehaan}
%===============================================================================

\textcite{dehaan2018power} address a widely acknowledged gap\sidenote{\textbf{De Haan \& Rotmans (2018)}---Agency-centered framework gives explanatory primacy to deliberate, value-driven action.}: the inadequate representation of actors and agency in transitions studies. They propose giving ``explanatory primacy to agency,'' presenting transitions as consequences of deliberate, value-driven actions rather than systemic forces.

Their framework introduces a four-role typology for transformative actors:

\marginnote{%
\scriptsize
\begin{tabular}{@{}p{1.8cm}p{2.2cm}@{}}
\toprule
\textbf{Role} & \textbf{Function} \\
\midrule
Frontrunners & Introduce alternatives \\
Connectors & Link actors \& solutions \\
Topplers & Change institutions \\
Supporters & Adopt and endorse \\
\bottomrule
\end{tabular}\\[0.3em]
Agency types and transition roles
}

Crucially, actors are not ``part of'' systems but \textit{affiliate} with them---holding multiple affiliations simultaneously. This framing enables analysis of boundary-spanning actors who connect different systems, and explains why the same organization might appear in multiple role positions across different contexts.

The framework also introduces \textbf{streams}---societal value sets enabled by available knowledge and organizing principles. Streams allow alliance formation around shared ideals rather than specific solutions. Actors form alliances of increasing scale: initiatives (frontrunner-driven), networks (connector-driven), and movements (toppler-driven).

However, the typology covers only \textit{transformative} actors---those maintaining the status quo are explicitly bracketed. This is a significant limitation, as transitions are constituted by conflict between incumbent and challenger coalitions. The framework also lacks systematic methods for empirical identification: how does one identify a ``frontrunner'' or ``toppler'' without circular reasoning from outcomes?


%===============================================================================
\section{Intermediary Functions and Typologies}[Intermediary Functions]
\labsec{background:intermediaries}
%===============================================================================

Given this dissertation's focus on intermediary organizations, we examine in depth how the literature has conceptualized intermediation---tracing its evolution from innovation studies into sustainability transitions research \cite{vanlente2020intermediaries}.\sidenote{\textbf{Conceptual evolution}---From bilateral brokers to systemic actors operating at network and system levels.}

\subsection{From Bilateral to Systemic Intermediaries}

The intellectual roots trace to innovation studies in the late 1990s \cite{howells2006intermediation}. The pivotal contribution came with \textcite{vanLente2003roles} on ``Roles of Systemic Intermediaries in Transition Processes,'' distinguishing \textbf{systemic intermediaries} from traditional bilateral brokers. They argued these actors function at the \emph{system or network level} rather than mediating one-to-one exchanges, identifying three core functions: \textbf{articulation} of options and demand, \textbf{alignment} of actors, and \textbf{support of learning processes}.

The intermediary concept evolved as it merged with sustainability transitions theory. Within the MLP, intermediaries came to be understood as actors positioned \emph{between} levels---connecting niche innovations to regime transformation. In Strategic Niche Management, intermediaries perform critical aggregation activities that construct the ``global niche'' as a cognitive and social reality.

\begin{insightbox}[The Aggregation Function]
Intermediaries do not merely connect actors---they \emph{aggregate} dispersed knowledge into coherent narratives. By circulating between local projects, they identify patterns, codify lessons, and construct shared understanding. Without this aggregation, individual experiments remain isolated and fail to accumulate into transition momentum.
\end{insightbox}

\subsection{The Kivimaa Typology}

The most influential recent contribution is \textcite{kivimaa2019innovation},\sidenote{\textbf{Kivimaa et al. (2019)}---Systematic review proposes five-type intermediary typology.} proposing a five-type typology:

\marginnote{%
\scriptsize
\begin{tabular}{@{}p{1.4cm}p{2.6cm}@{}}
\toprule
\textbf{Type} & \textbf{Key features} \\
\midrule
Systemic & Across levels; transition agenda \\
Regime & Reformer within regime \\
Niche & Advances innovations \\
Process & Neutral broker \\
User & Connects users to tech \\
\bottomrule
\end{tabular}\\[0.3em]
The Kivimaa intermediary typology
}

\textbf{Systemic intermediaries} operate across niche, regime, and landscape levels with explicit transition agendas. National agencies like Finland's Sitra exemplify this type. \textbf{Regime-based intermediaries} remain associated with prevailing regimes but promote transition---reformers rather than revolutionaries. \textbf{Niche intermediaries} work to advance specific innovations, developing voice and formalization. \textbf{Process intermediaries} facilitate change without explicit system-transformation agendas, earning trust through neutrality. \textbf{User intermediaries} connect users to new technologies and communicate preferences to developers.

Yet empirical studies reveal\sidenote{\textbf{Porous boundaries}---Intermediaries often span types and shift strategically.} that boundaries between types are porous---intermediaries often exhibit characteristics of multiple types simultaneously and may shift profiles strategically over time.

\subsection{Function Clusters}

The extensive documentation of intermediary functions reveals four clusters:

\textbf{Knowledge functions}: gathering, processing, and generating knowledge; brokering and transferring across domains; disseminating information; facilitating learning; aggregating lessons across experiments.

\textbf{Network functions}: creating and managing networks; configuring and aligning interests; building partnerships and coalitions; connecting actors across boundaries; building relational capital---trust, shared understanding, collaborative capacity.

\textbf{Vision and expectation work}: articulating futures; sense-making through narrative construction; framing problems and solutions in ways that mobilize action; managing ``promises and requirements'' dynamics.

\textbf{Institutional functions}: advocating for policy change; legitimation work to build social acceptance; shaping regulatory frameworks; supporting institutional infrastructure development.


\subsection{Related Concepts}

The intermediary concept operates alongside several related terms. \textbf{Boundary spanners} emphasize individual competencies for bridging knowledge domains, particularly at the science-policy interface \cite{tushman1977boundary}. \textbf{System builders} engage in strategic action to create favorable conditions, more directive than facilitative. \textbf{Transition brokers} represent systemic intermediaries fulfilling orchestration functions. \textbf{Orchestrators} exercise leadership and make strategic decisions, contrasting with facilitative orientations.


%===============================================================================
\section{Kanda et al.: Intermediary Dynamics and System Levels}[Kanda et al.]
\labsec{background:dynamics}
%===============================================================================

\textcite{kanda2020conceptualising} address a specific gap\sidenote{\textbf{Kanda et al. (2020)}---Emphasizes temporal dynamics and ecology perspectives.}: how to conceptualize intermediary activities \textit{beyond} individual projects at the system level. Their key move: shift from \textit{systemic intermediaries} (actor category) to \textit{systemic intermediation} (activity description). Any intermediary may perform systemic activities; dedicated ``systemic intermediaries'' may also perform non-systemic activities.

They propose four levels of intermediation based on the relational criterion ``in-between what?'': \textbf{Level 0} (bilateral exchanges between individual entities), \textbf{Level 1} (mediation within a single network), \textbf{Level 2} (spanning different network types), and \textbf{Level 3} (engagement between actors, networks, and institutions). Activities range from one-to-one consulting at Level 0 to lobbying and policy translation at Level 3.

Key empirical finding: \textbf{non-systemic activities dominate}.\sidenote{\textbf{Phase-dependent roles}---Niche intermediaries dominate early; systemic intermediaries peak during acceleration; diffusion intermediaries become critical for mainstreaming.} Most intermediation is bilateral despite the literature's focus on ``systemic intermediaries.'' Systemic activities require financial stability, longevity, and sometimes neutrality---they are resource-intensive. Mandate matters as much as resources: limited-resource intermediaries with institutional change mandates can operate at Level 3, while well-resourced intermediaries with bilateral mandates may concentrate on Level 0.

\subsection{Emergence Pathways}

Three intermediary emergence pathways shape characteristics and capabilities:

\textbf{Purpose-built}: specifically established to intermediate, typically by governments creating agencies with explicit transition mandates. These begin with formal mandates but must build legitimacy and networks.

\textbf{Evolved}: existing organizations that adopt intermediary roles as conditions change. An industry association may evolve from representing incumbent interests to facilitating transitions. These leverage existing relationships but may carry institutional baggage.

\textbf{Emergent}: arising spontaneously in response to system gaps, often without initially recognizing their intermediary function. Community energy groups and informal networks may gradually assume broader roles. These possess grassroots legitimacy but may lack resources.

\subsection{Evolution Across Transition Phases}

What intermediaries need to do---and which types matter most---changes as transitions progress.

In \textbf{pre-development}, intermediaries support early experimentation, build initial networks among pioneers, and facilitate learning from failures. Niche and process intermediaries are particularly important.

During \textbf{take-off}, functions shift toward knowledge aggregation, building broader coalitions, and pooling resources for scaling. Systemic intermediaries become more important.

\textbf{Acceleration} demands stronger institutional support, capacity building for mainstream actors, and substantial resource mobilization. ``Diffusion intermediaries'' help adopters navigate new practices.

In \textbf{stabilization}, some intermediaries cease to exist as roles become unnecessary---a sign of success. Others find new opportunities or pivot to remaining niches.

A crucial finding: ``without intermediary influence, pre-development is unlikely to lead to acceleration.'' Intermediaries provide continuity across phases, carrying lessons and relationships even as specific projects end.

\subsection{Ecologies of Intermediaries}

Perhaps the most significant recent advance is recognizing that transitions involve \emph{multiple} intermediaries whose interactions shape outcomes.

\begin{roledef}[Ecology of Intermediaries]
An \textbf{ecology of intermediaries} comprises a variety of intermediary actors with different missions, mandates, and levels of agency that collectively connect actors and resources at different scales within a transition arena.
\end{roledef}

Within ecologies, intermediaries exhibit both \textbf{complementary and competitive dynamics}.\sidenote{\textbf{Beyond single intermediaries}---Most research examines individual intermediaries. The ecology concept redirects attention to how multiple intermediaries interact, complement, and compete.} Complementarity arises when different types address different functions: systemic intermediaries articulating visions while process intermediaries manage projects. Yet intermediaries also \textbf{compete} for resources, attention, legitimacy, and influence.

\textbf{Ecology gaps}---missing intermediary functions---can hinder transition progress. A transition arena may have strong niche intermediaries but lack systemic coordination, or possess visionary leadership but insufficient process support. Mapping these gaps is essential for transition governance but requires systematic methods currently lacking.


%===============================================================================
\section{Geographic and Methodological Considerations}[Geography \& Methods]
\labsec{background:context}
%===============================================================================

\textbf{Geographic context} matters substantially.\sidenote{\textbf{Context matters}---Geographic variation and methodological constraints shape knowledge production.} While European research dominates---particularly from the Netherlands, Finland, Germany, and UK---patterns vary across governance systems. The Sustainability Transitions Research Network's ``Transitions in Global South'' initiative has catalyzed research beyond Northern contexts, revealing that intermediary roles and ecosystem configurations differ with institutional environments, development trajectories, and power structures.

\subsection{The Dominance of Designed Data}

A striking but underappreciated feature of role constellation research is its near-total reliance on \textbf{designed data}---data created specifically for research purposes.\sidenote{\textbf{Designed vs. organic data}---Groves (2011) distinguishes data created \emph{for} research (surveys, interviews) from data existing \emph{independently} of research (digital traces, transaction logs).} Interviews, surveys, focus groups, workshop protocols---these are the dominant data sources in transitions research. \textcite{zolfagharian2019studying} found that 82\% of sustainability transitions studies use qualitative methods, with 58\% relying on interviews as the primary data source. Only 9\% employ quantitative methods, and most of these use survey instruments designed for specific studies.

This reliance on designed data has consequences. Interview-based research captures what actors \emph{say} about their roles---their self-perceptions, interpretations, and narratives. Surveys impose researcher-defined categories that respondents must fit themselves into. Both approaches filter role information through actors' self-presentation and researchers' conceptual frameworks.

More fundamentally, designed data does not scale. Each interview requires researcher time; each survey requires respondent cooperation. A study of thirty intermediaries represents substantial investment. A study of three hundred would be prohibitively expensive. A study of three thousand is essentially impossible with designed data methods.

\subsection{Organic Data: The Uncharted Territory}

\textbf{Organic data}---data generated independently of research purposes---offers an alternative.\sidenote{\textbf{Organic data sources}---Websites, social media, news coverage, funding databases, visual communications---all generated for organizational purposes, not research consumption.} Organizations create websites to present themselves to stakeholders. They post on social media to engage audiences. They appear in news coverage whether they seek it or not. They receive and make investments documented in financial databases. They produce visual communications for marketing and branding. None of this is created for researchers---but all of it carries information about organizational roles.

The potential is substantial. Where designed data captures what actors \emph{say} they do, organic data can reveal what they \emph{actually} do: whom they connect with, how they present themselves, how others perceive them, where money flows. Where designed data reaches dozens of organizations, organic data can cover thousands. Where designed data offers snapshots, organic data can be collected longitudinally.

Yet organic data remains largely uncharted territory in transitions research. \textcite{dehaan2019datadriven} observe that ``not only is there no such database but moreover it seems contrary to the way data are used in transitions research.'' The field has developed sophisticated theories about role dynamics but has not developed the data infrastructure to test them at scale.

This is not merely a technical gap. It reflects deeper assumptions about what counts as valid evidence. Designed data privileges actor voice---the interpretations and meanings actors themselves assign to their roles. Organic data privileges behavioral traces---what actors do regardless of what they say. Both have validity; neither is sufficient. The question is how to combine them.

\subsection{Methodological Trade-offs}

Existing methodological approaches face persistent trade-offs.\sidenote{\textbf{Methodological tensions}---Depth vs. scale; meaning vs. pattern; theory vs. data.} Qualitative case studies provide rich understanding of how intermediaries navigate complex contexts but cannot identify population-level patterns. Social network analysis maps relational structures but typically relies on single data sources and struggles to capture meaning. Computational and text-analytic approaches offer scale but often begin with available data rather than theoretical questions.

Each approach illuminates some aspects while obscuring others. No single method captures the full terrain of intermediary role constellations. The challenge is integration: how to combine designed data's depth with organic data's scale, qualitative richness with quantitative pattern detection, theoretical guidance with data-driven discovery.


%===============================================================================
\section{Synthesis: Convergences, Divergences, and Gaps}[Synthesis]
\labsec{background:synthesis}
%===============================================================================

Having reviewed individual contributions, we now synthesize across them to identify what these frameworks share, where they diverge, and what gaps remain.


\subsection{Points of Convergence}

Despite originating from different intellectual traditions,\sidenote{\textbf{Five convergences}---Despite different theoretical origins, frameworks agree on key principles.} the four frameworks converge on several key insights:

\textbf{Actor-process distinction}: All four distinguish between actors-as-entities and the processes they perform. De Haan \& Rotmans separate actors from transformative change; Wittmayer et al. separate actors from roles; Kanda et al. separate intermediaries from intermediation; Fischer \& Newig separate actors from agency.

\textbf{Multiple positions}: Actors are not bound to single categories. De Haan \& Rotmans emphasize multiple affiliations to systems; Wittmayer et al. note actors occupy multiple roles; Kanda et al. show intermediaries operate across levels; Fischer \& Newig find actors belong to different categories simultaneously.

\textbf{Temporal instability}: Roles and positions change over transition phases. De Haan \& Rotmans trace role trajectories; Wittmayer et al. describe roles as temporary stabilizations; Fischer \& Newig find roles ``erratic.''

\textbf{Relational over categorical}: Position relative to others matters more than intrinsic attributes. De Haan \& Rotmans emphasize affiliation structure; Wittmayer et al. focus on role constellations; Kanda et al. apply an ``in-between what?'' criterion; Fischer \& Newig highlight dependencies between actors.

\textbf{Typology limitations}: All find existing categories inadequate. De Haan \& Rotmans argue MLP cannot capture actor dynamics; Wittmayer et al. note transitions lack role vocabulary; Kanda et al. find ``systemic intermediary'' underspecified; Fischer \& Newig identify four unintegrated typologies coexisting.

\subsection{Points of Divergence and Shared Gaps}

While these frameworks converge on key principles, they diverge on agency sources (values vs. social structure vs. mandates), change mechanisms (strategic action vs. role negotiation vs. higher-level activities), and empirical scales (from neighborhoods to fields). More striking than these divergences are the gaps all frameworks share.\sidenote{\textbf{Six gaps}---These limitations motivate the cartographic approach developed in this dissertation.}

None offers systematic \textbf{operationalization}. The frameworks provide concepts---frontrunners, connectors, role constellations---but not methods for identifying them empirically. How do you find a frontrunner without already knowing who drove the transition? Researcher interpretation fills the void, but interpretation varies.

None enables \textbf{population-level analysis}. Each works at case level. How many intermediaries exist in a given ecosystem? What is the distribution across types? Do certain configurations cluster geographically? These questions cannot be answered with existing approaches.

None provides \textbf{cross-modal validation}. Organizations claim roles on their websites. Do they actually perform those roles? No framework systematically compares self-description with observed behavior. The assumption seems to be that what actors say matches what they do.

None specifies \textbf{boundaries} clearly. Role constellations, networks, systems---all invoked, but where does one end and another begin? Membership criteria remain implicit. This makes comparison across studies difficult.

None develops systematic \textbf{power analysis}. Dependencies, resources, constellations---all gesture toward power without analyzing it directly \cite{avelino2024just}. Who dominates? Whose roles constrain others? What asymmetries shape constellation dynamics?

None addresses \textbf{failure and absence}. What happens when critical roles go unfilled? When intermediaries exit? When constellations contain gaps? The literature focuses on presence and function; absence remains theorized but unmapped.

None conceptualizes \textbf{phase-out and exnovation roles}. The X-curve framework \sidecite{hebinck2022xcurve} demonstrates that sustainability transitions require not only build-up dynamics (experimentation, acceleration, institutionalization) but equally breakdown dynamics (destabilization, chaos, phase-out). Yet the role typologies reviewed above---frontrunners, connectors, intermediaries, network builders---are exclusively oriented toward the build-up side of the X-curve.\sidenote{\textbf{The X-curve blind spot}---Where are the destabilizers? The phase-out managers? The regime dismantlers? Existing typologies cannot answer because they do not ask.} Who facilitates managed decline? Who coordinates the phase-out of fossil fuel infrastructure? Who supports workers and communities through regime breakdown? These questions remain unasked because the conceptual vocabulary does not exist. This represents perhaps the most significant gap: half of transition dynamics---the breakdown half---lacks corresponding role conceptualization.

\subsection{Toward Integration}

The governance dimension adds a fourth consideration. \textcite{loorbach2017transitions} distinguish three approaches to governance in transitions: \textit{analyzing} governance (understanding how societal agency interacts with institutional change), \textit{evaluating} governance (assessing whether transition-oriented policies work), and \textit{experimental} governance (intervening through arenas, scenarios, and experiments).\sidenote{\textbf{Three governance approaches}---Analysis, evaluation, and experimentation \cite{loorbach2017transitions}.} Role constellations matter for all three: analysis requires mapping who performs what functions; evaluation requires understanding how role configurations enable or constrain policy effectiveness; experimentation requires identifying intervention points where new roles can be created or existing ones reconfigured.

Despite their differences, these frameworks converge on a shared insight: understanding role constellations requires attending to three dimensions simultaneously.\sidenote{\textbf{Three dimensions}---Function, relation, time. Existing typologies emphasize one while backgrounding others.}

The \textbf{functional dimension} asks what activities are performed. Kivimaa's intermediary functions, Kanda et al.'s systemic activities, De Haan \& Rotmans' frontrunner/connector/toppler roles---all characterize what actors \emph{do}. But function alone is insufficient. The same activity---brokering knowledge, for instance---means different things depending on who performs it and between whom.

The \textbf{relational dimension} asks where action occurs. Kanda et al.'s ``in-between what?'' criterion, Wittmayer et al.'s constellation concept, Fischer \& Newig's actor dependencies---all recognize that roles exist only relative to other roles. An intermediary is not an intermediary in isolation; the role emerges from the space between other actors. Position in relational space determines what functions are available and effective.

The \textbf{temporal dimension} asks how positions change. De Haan \& Rotmans trace role trajectories across transition phases; Wittmayer et al. describe roles as temporary stabilizations that shift as transitions unfold; Fischer \& Newig find that roles are ``erratic''---actors move between categories as circumstances change. Static typologies miss this dynamism.

A comprehensive framework would characterize positions along all three dimensions---specifying \textit{what} actors do, \textit{where} they are positioned relationally, and \textit{how} this changes over time. The challenge is methodology: how do you map all three simultaneously at scale?


%===============================================================================
\section{Seven Persistent Problems---and a Cartographic Response}[Seven Problems]
\labsec{background:problems}
%===============================================================================

Building on this synthesis, we identify seven persistent problems that limit current understanding of role constellations.\sidenote{\textbf{Problems as opportunities}---Each problem suggests what better cartography would accomplish, not solving problems definitively but reframing them productively.} These are not merely gaps to be filled but \textbf{productive tensions} pointing toward cartographic opportunities.

\begin{enumerate}[leftmargin=*, itemsep=0.8em]

\item \textbf{Definitional ambiguity.} ``Role'' variously denotes position, function, identity, or behavior---conflations that impede cumulation.

\emph{Cartographic reframing}: Rather than seeking definitional consensus, cartography \textbf{distinguishes layers}. Position can be mapped relationally (where actors sit in network space); function can be mapped through activity traces (what actors do); identity can be mapped through self-presentation. Different layers can be \textbf{superimposed} to reveal how position, function, and identity align or diverge.

\item \textbf{Static categorization.} Typologies present roles as fixed boxes, yet empirical reality shows fluidity---actors shifting roles, occupying multiple roles, varying by context.

\emph{Cartographic reframing}: Categories are \textbf{snapshots}; cartography traces \textbf{trajectories}. Rather than asking ``which box does this intermediary belong in?'' we ask ``where has this intermediary been, where is it now, and in what direction is it moving?''

\item \textbf{Actor-centric bias.} Roles are typically treated as actor attributes rather than relational accomplishments, obscuring how roles emerge through interaction.

\emph{Cartographic reframing}: Cartography is inherently \textbf{relational}---positions exist only relative to other positions; territories are defined by boundaries with neighbors; paths connect locations. Mapping intermediaries means mapping the spaces \emph{between} them, the connections that constitute positions.

\item \textbf{Normative loading.} Role concepts carry implicit judgments---intermediaries good, incumbents bad; bridging valuable, blocking problematic.

\emph{Cartographic reframing}: Maps can be \textbf{descriptive}---charting terrain without evaluating it. A map shows where mountains are without claiming mountains are good. Cartography enables normative analysis while maintaining descriptive foundations.

\item \textbf{Methodological constraints.} Small-N qualitative studies provide depth but cannot identify population-level patterns; computational approaches offer scale but lack theoretical meaning.

\emph{Cartographic reframing}: Multimodal cartography \textbf{integrates} rather than chooses. Theory guides what features to map; multiple data sources triangulate positions; qualitative insight interprets quantitative patterns.

\item \textbf{Theoretical fragmentation.} Different frameworks emphasize different roles without integration. MLP, TIS, SNM, and transition management each offer partial maps that don't align.

\emph{Cartographic reframing}: Fragmentation reflects \textbf{different mapping conventions}. Rather than seeking a single master framework, cartography can \textbf{translate between} frameworks, showing how different conceptual schemes relate.

\item \textbf{Build-up bias.} All role typologies reviewed above---frontrunners, connectors, intermediaries, brokers---focus on the innovation and acceleration side of transitions. The breakdown dynamics captured in the X-curve \parencite{hebinck2022xcurve}---destabilization, phase-out, managed decline---lack corresponding role concepts.\sidenote{\textbf{Half the X-curve is unmapped}---If transitions require both build-up and breakdown, existing frameworks map only one arm.}

\emph{Cartographic reframing}: Cartography can \textbf{reveal absences}. By systematically mapping what actors \emph{do} claim and perform, we can identify what functions remain \emph{unclaimed}---including the phase-out roles that transition theory suggests should exist but our empirical mapping may fail to find.

\end{enumerate}

These reframings do not dissolve the problems. Definitional ambiguity persists; static categorization remains a temptation; actor-centric bias is hard to shake. But cartography offers different ways to work with these tensions. Instead of seeking the ``right'' definition of intermediary, we can map how different definitions highlight different features of the same actors. Instead of building better typologies, we can trace how actors move through role space over time. The problems become generative rather than paralyzing.


%===============================================================================
\section{Chapter Summary}
\labsec{background:summary}
%===============================================================================

\begin{summarybox}
This chapter has surveyed how sustainability transitions research conceptualizes roles, with particular attention to intermediaries as the focal actors of this dissertation:

\begin{itemize}[leftmargin=*, itemsep=0.3em]
\item \textbf{Role constellations} provide a meso-level lens for analyzing what functions actors perform and how these combine---distinct from actor constellations that merely catalog who is present.

\item Roles are \textbf{malleable and relational}: they emerge through interaction, shift across transition phases, and gain meaning through position within constellations.

\item \textbf{Key contributions} reviewed include: Fischer \& Newig's (2016) systematic typology mapping; Wittmayer et al.'s (2017) introduction of role theory; De Haan \& Rotmans' (2018) agency-centered framework; and Kivimaa et al.'s (2019) and Kanda et al.'s (2020) intermediary syntheses.

\item \textbf{Intermediaries} have evolved from bilateral brokers to systemic actors operating across niche-regime-landscape levels. The Kivimaa typology distinguishes systemic, regime-based, niche, process, and user intermediaries---but boundaries are porous.

\item Intermediaries perform four function clusters---\textbf{knowledge work, network building, vision articulation, and institutional change}.

\item \textbf{Ecologies of intermediaries} shape transitions through complementary and competitive dynamics. Ecology gaps can stall progress.

\item \textbf{Seven persistent problems}---definitional ambiguity, static categorization, actor-centric bias, normative loading, methodological constraints, theoretical fragmentation, and build-up bias---limit current understanding but point toward cartographic opportunities.

\item The field's core limitation is \textbf{categorical thinking} where \textbf{cartographic thinking} is needed.
\end{itemize}

\noindent These seven problems directly motivate our first sub-research question: \textit{What conceptual framework enables theory-guided but empirically open analysis of role constellations?} The literature provides conceptual vocabulary---role constellations, intermediary functions, ecologies---but lacks operationalization at scale. The cartographic reframings outlined above constitute this dissertation's theoretical response.
\end{summarybox}

\vspace{1em}

The literature provides rich conceptual resources---typologies, function inventories, phase frameworks, ecology concepts. What it lacks is methodology adequate to its own insights. Scholars have developed sophisticated ways to \emph{think} about role constellations. Tools to \emph{map} them systematically remain scarce. Theory outpaces method.

This gap is where the dissertation intervenes. The next chapter presents the Design Science Research framework\sidenote{\textit{Next:} Chapter~\ref{ch:method}---Design Science Methodology: how artifacts and infrastructure enable systematic multimodal cartography.} that guides artifact construction, detailing Meta-Design Requirements and data collection workflows across six modalities. The EIT ecosystem introduced earlier provides the empirical laboratory; the methodology chapter provides the tools to map it.

\end{document}
